\lecture{10}{26 March.}{}
Let \(X_i\) be the random variable representing the earning of option \(i\) after a year. Then 
\begin{align*}
    \mathbb{E} [X_1] &= \$ 1 \cdot 1 = \$ 1. \\
    \mathbb{E} [X_2] &= \$ 1000 \cdot \frac{1}{1000} + \$ 0 \cdot \frac{999}{1000} = \$ 1. \\
    \mathbb{E} [X_3] &= \$ \frac{2000}{999} \cdot \frac{999}{1000} - \$ 1000 \cdot \frac{1}{1000} = \$ 1.
\end{align*} 

Thus, three options all have same expectation, which means expectation can not reflect all. Hence, we need variance, which can measure how far a random typically is from its expectation.

\begin{definition}[Variance]
    Let \(X\) be a random variable on a probability space \((S, \mathbb{P} )\), and let \(\mu = \mathbb{E} [X]\). We define the variance to be 
    \[
        \Var_{}(X) = \mathbb{E} \left[ (X - \mu )^2 \right]. 
    \]   
\end{definition}

Thus, 
\begin{align*}
    \Var (X_1) &= \mathbb{E} [(X_1 - 1)^2] = (1 - 1)^2 \cdot 1 = 0 \\
    \Var [X_2] &= \mathbb{E} [(X_2 - 1)^2] = (1000 - 1)^2 \cdot \frac{1}{1000} + (0 - 1)^2 \cdot \frac{999}{1000} \thickapprox 999. \\
    \Var [X_3] &= \mathbb{E} [(X_3 - 1)^2] = \left( \frac{2000}{999} - 1 \right)^2 \cdot \frac{999}{1000} + (-1000 - 1)^2 \cdot \frac{1}{1000} \thickapprox 1003. 
\end{align*}

\begin{claim}
    For any random variable, 
    \[
        \Var(X) = \mathbb{E} [X^2] - \mathbb{E} [X]^2.
    \]
\end{claim}

\begin{explanation}
    \begin{align*}
        \Var (X) &= \mathbb{E} [(X - \mu )^2] = \mathbb{E} [X^2 - 2 \mu X + \mu ^2] \\
        &= \mathbb{E} [X^2] - 2\mu \mathbb{E} [X] + \mu ^2 \\
        &= \mathbb{E} [X^2] - \mathbb{E} [X]^2 
    \end{align*}
    since \(\mu = \mathbb{E} [X]\). 
\end{explanation}

\begin{corollary}
    \(\mathbb{E} [X^2] \ge \mathbb{E} [X]^2\). 
\end{corollary}
\begin{proof}
    \[
        0 \le \mathbb{E} [(X - \mu )^2] = \Var(X) = \mathbb{E} [X^2] - \mathbb{E} [X]^2.
    \]
\end{proof}

\begin{remark}
    Sometimes it is easier to compute 
    \[
        \mathbb{E} [X(X - 1)] = \mathbb{E} [X^2] - \mathbb{E} [X],
    \]
    then 
    \[
        \Var(X) = \mathbb{E} [X(X-1)] + \mathbb{E} [X] - \mathbb{E} [X]^2.
    \]
\end{remark}

\begin{definition}
    The standard deviation of a random variable \(X\), denoted \(\sigma (X)\), is 
    \[
        \sigma (X) = \sqrt{\Var(X)}. 
    \]  
    We sometimes denote to variance by \(\sigma ^2(X)\). 
\end{definition}

\begin{proposition}[Markov's Inequality]
    Let \(X \ge 0\) be a non-negative random variable on a probability space \((S, \mathbb{P} )\), and let \(a > 0\) be a positive real number. Then, 
    \[
        \mathbb{P} (X \ge a) \le \frac{\mathbb{E} [X]}{a}.
    \]   
\end{proposition}

\begin{proof}
    \begin{align*}
        \mathbb{E} [X] &= \sum_{x} x \mathbb{P} (X = x) = \underbrace{\sum_{x: x < a} x\mathbb{P} (X = x)}_{\ge 0} + \sum_{x : x \ge a} x \mathbb{P} (X = x) \\
        &\ge \sum_{x: x \ge a} a \mathbb{P} (X = x) = a \sum_{x: x \ge a} \mathbb{P} (X = x)  = a \mathbb{P} (X \ge a).    
    \end{align*}
\end{proof}

\begin{remark}
    Markov's inequality says that small expectation implies random variable is small with high probability, but Markov's inequality does not say if expectation is large, then  random variable large with high probability. In fact, this is not always true: Consider 
    \[
        X_n = \begin{dcases}
            n^2, &\text{ with probability } \frac{1}{n} ;\\
            0, &\text{ with probability } 1 - \frac{1}{n} .
        \end{dcases}
    \]
    Then,
    \[
        \mathbb{E} [X_n] = n^2 \cdot \frac{1}{n} = n \to \infty \text{ as } n \to \infty. 
    \]
    However, 
    \[
        \mathbb{P} (X_n \neq 0) = \frac{1}{n} \to 0 \text{ as } n \to \infty. 
    \]
\end{remark}

\begin{theorem}[Chebyshev's inequality]
    Let \(X\) be a random variable with \(\mu = \mathbb{E} [X]\) on a probability space \((S, \mathbb{P} )\), and let \(a > 0\) be a real number. Then,
    \[
        \mathbb{P} (\vert X - \mu \vert \ge a) \le \frac{\Var(X)}{a^2}.
    \]    
\end{theorem}

\begin{proof}
    We can apply Markov's inequality to \(Y = (X - \mu)^2 \ge 0\). Thus,
    \[
        \mathbb{P} (Y \ge a^2) \le \frac{\mathbb{E} [Y]}{a^2} = \frac{\mathbb{E} [(X - \mu )^2]}{a^2} = \frac{\Var(X)}{a^2}.
    \] 
\end{proof}

\begin{remark}
    Chebyshev's inequality implies small variance implies \(X\) is close to its expectation with high probability. 
\end{remark}

\begin{corollary}
    The probability that \(X\) is more than \(k\) standard deviation from its expectation is \(\le \frac{1}{k^2}\).   
\end{corollary}
\begin{proof}
    \begin{align*}
        \mathbb{P} (\vert X - \mu  \vert \ge k \sigma (X) ) \le \frac{\Var(X)}{(k \sigma (X))^2} = \frac{\Var(X)}{k^2 \sigma ^2(X)} = \frac{1}{k^2}.
    \end{align*}
\end{proof}

\begin{observation}
    Let \(X\) be a random variable, \(a, b \in \mathbb{R}\). Then, 
    \[
        \mathbb{E} [aX + b] = a \mathbb{E} [X] + b
    \] by linearity, and 
    \begin{align*}
        \Var(aX + b) &= \mathbb{E} \left[ ((aX+b) - \mathbb{E} [aX+b])^2 \right] = \mathbb{E} \left[ ((aX + b) - (a\mathbb{E} [X] + b))^2 \right] \\
        &= \mathbb{E} \left[ a^2 \left( X - \mathbb{E} [X] \right)^2  \right] = a^2 \mathbb{E} \left[ (X - \mathbb{E} [X])^2 \right] = a^2 \Var(X).  
    \end{align*} 
\end{observation}

\section{Discrete Distribution}
\begin{definition}[Bernoulli distribution]
    If \(X \sim \mathrm{Ber}(p) \), i.e. \(X\) has the \(\mathrm{Ber}(p) \) distribution, then 
    \[
        X = \begin{dcases}
            1, &\text{ with probability } p ;\\
            0, &\text{ with probability } 1-p .
        \end{dcases}
    \]
    Then, 
    \[
        \mathbb{E} [X] = 1 \cdot p + 0 \cdot (1 - p) = p
    \]    
    and 
    \[
        \Var(X) = \mathbb{E} [X^2] - \mathbb{E} [X]^2 = \mathbb{E} [X] - \mathbb{E} [X]^2 = p - p^2 = p(1 - p).
    \]
\end{definition}
   
\begin{definition}[Binomial distribution]
    If we run \(n\) independent trials of a random experiment of which succeeds with probability \(p\). We count the number of successes, where the possible values are \(0, 1, 2, \dots , n - 1, n\). Thus, we can define \(\mathrm{Bin}(n, p), n \in \mathbb{N} , 0 \le p \le 1 \) to be the number of successes in the random experiment. 
\end{definition}

\begin{eg}
    \hfill 
    \begin{itemize}
        \item \(\#\) of heads in \(n\) coin tosses. 
        \item \(\#\) of odd rolls in \(n\) rolls of a die. 
        \item \(\#\) of errors when replicating DNA. 
        \item \(\#\) of supporters of a particular party in a poll.      
    \end{itemize}
\end{eg}

The probability mass function of \(\mathrm{Bin}(n, p)\) is 
\[
    \mathbb{P} (X = k) = \binom{n}{k} p^k (1 - p)^{n - k}.
\] 

Hence,
\[
    \mathbb{E} [X] = \sum_{k=0}^n k \cdot \mathbb{P} (X = k) = \sum_{k=0}^n k \binom{n}{k} p^k (1 - p)^{n-k}.  
\]
Alternatively,
\[
    X = \sum_{i=1}^n X_i, \text{ where } X_i = \begin{dcases}
        1, &\text{ if } i \text{-th trial is successful}  ;\\
        0, &\text{ if not}  ,
    \end{dcases} 
\]
and each \(X_i \sim \mathrm{Ber}(p) \). Thus,
\begin{align*}
    \mathbb{E} [X] &= \mathbb{E} \left[ \sum_{i=1}^n X_i  \right] = \sum_{i=1}^n \mathbb{E} [X_i] = \sum_{i=1}^n p = np.   
\end{align*} 

Recall the Binomial Theorem
\[
    (x + y)^n = \sum_{k=0}^n \binom{n}{k} x^k y^{n-k}. 
\]
Thus,
\begin{align*}
    n(x + y)^{n - 1} &= \frac{\mathrm{d}}{\mathrm{d}x} (x + y)^n = \frac{\mathrm{d}}{\mathrm{d}x} \left[ \sum_{k=0}^n \binom{n}{k} x^k y^{n-k} \right] = \sum_{k=0}^n k \binom{n}{k} x^{k-1} y^{n - k}.  
\end{align*}
Multiplying \(x\) on both sides, then 
\[
    nx(x + y)^{n - 1} = \sum_{k=0}^n k \binom{n}{k} x^k y^{n - k}.  
\] 
Now substituting \(x = p\) and \(y = 1- p\), then 
\[
    np = np (p + (1 - p))^{n - 1} = \sum_{k=0}^n k \binom{n}{k} p^k (1 - p)^{n - k} = \mathbb{E} [X]. 
\]  

Now we compute \(\Var(X)\). Note that 
\[
    \Var(X) = \mathbb{E} [X^2] - \mathbb{E} [X]^2
\] 
and 
\[
    \mathbb{E} [X^2] = \sum_{k=0}^n k^2 \binom{n}{k} p^k (1 - p)^{n - k}. 
\]
Differentiable with respect to \(x\), then we have 
\begin{align*}
    \frac{\mathrm{d}}{\mathrm{d}x} \left[ nx(x + y)^{n - 1} \right] &= \frac{\mathrm{d}}{\mathrm{d}x} \left[ \sum_{k=0}^n k \binom{n}{k} x^k y^{n - k}  \right].    
\end{align*} 
Hence,
\[
    n(x + y)^{n - 1} + n(n - 1)x(x + y)^{n - 2} = \sum_{k=0}^n k^2 \binom{n}{k} x^{k - 1} y^{n - k}. 
\]
Multiplying by \(x\), we have 
\[
   nx(x + y)^{n - 1} + n(n - 1)x^2(x + y)^{n - 2} = \sum_{k=0}^n k^2 \binom{n}{k} x^{k} y^{n - k}, 
\] 
and substituting \(x = p, y = 1 - p\), we have 
\[
    np + n(n - 1)p^2 = \sum_{k=0}^n k^2 \binom{n}{k} p^k (1 - p)^{n - k} = \mathbb{E} [X^2]. 
\] 
Thus,
\begin{align*}
     \Var(X) &= \mathbb{E} [X^2] - \mathbb{E} [X]^2 = np + n(n - 1)p^2 - (np)^2 \\
     &= np - np^2 = np(1 - p) = n \cdot \Var(\mathrm{Ber}(p)).
\end{align*}

\begin{definition}[Poisson distribution]
    If there are rare independent events, and we count how many occurences in a given time period, then we denote it by \(\mathrm{Poi}(\lambda ) \), where \(\lambda \ge 0\) is a real number.  
\end{definition}

\begin{eg}
    \hfill
    \begin{itemize}
        \item earthquakes. 
        \item \(\#\) of customers eating store. 
        \item \(\#\) of phone calls.  
    \end{itemize}
    \(\lambda =\) expected \(\#\) of occurences.  
\end{eg}

Thus, possible values of \(X \sim \mathrm{Poi}(\lambda ) \) is \(0, 1, 2, \dots \).  